% Capitulo: Com que frequência os dados precisam atualizar?
\chapter{Com que frequência os dados precisam atualizar?}
\label{cap:frequencia}

Poucas decisões num projeto de dados são tomadas tão no automático quanto a frequência de
atualização. O stakeholder pede \enquote{o mais atualizado possível}, o time configura o
pipeline para rodar de hora em hora porque a ferramenta permite, e ninguém volta ao assunto.
Meses depois, descobre-se que o painel é aberto uma vez por semana, na reunião de
segunda-feira, e que todas as outras atualizações consumiram recurso sem que ninguém olhasse
o resultado.

A regra que este capítulo defende é simples de enunciar: atualizar o mais rápido que as fontes
permitem, mas só o quanto a decisão de negócio precisa. Aplicá-la exige entender como o
negócio decide, conhecer os limites das fontes, dividir a responsabilidade entre os papéis
do time e, principalmente, medir o uso dos próprios produtos de dados, porque a primeira
definição quase sempre está errada.

\section{O ritmo da decisão}
\label{sec:ritmo-decisao}

A pergunta \enquote{com que frequência você quer que o dado atualize?} tem uma resposta
previsível: o mais rápido possível. Ninguém pede dado velho. A pergunta certa é outra:
\enquote{quando e como você toma a decisão com esse número?}. A resposta a essa pergunta define
a frequência necessária.

Decisões operacionais, como despachar uma entrega, repor um item que acabou ou bloquear uma
transação suspeita, acontecem em minutos ou horas, e pedem dados atualizados no mesmo ritmo.
Decisões táticas, como ajustar a escala da semana ou revisar o preço de uma linha de produtos,
convivem bem com dados do dia anterior. Decisões estratégicas, como definir metas, aprovar
orçamento ou planejar uma safra, são tomadas uma vez por mês ou por ano, com dados fechados.

A Figura~\ref{fig:frequencia-consumo} cruza a frequência de atualização com a frequência com
que alguém realmente olha os números. Nos dois quadrantes alinhados, a atualização acompanha o
uso. Nos outros dois, há um problema. No quadrante do desperdício, o dado atualiza muito mais do
que é consultado: recurso gasto sem que ninguém aproveite. No quadrante do risco, muita gente
decide todo dia com um número que só atualiza de vez em quando, e as decisões são tomadas com
dado velho sem que ninguém perceba.

\begin{figure}[htbp]
  \centering
  % Frequencia de atualizacao contra frequencia de consumo (cap. 16).
\begin{tikzpicture}[diagrama, x=1mm, y=1mm]
  % quadrantes (os da diagonal, alinhados, em branco; os desalinhados, em cinza)
  \node[caixa, fill=fundo, minimum width=46mm, minimum height=26mm, text width=40mm] at (25,43)
    {\textbf{Desperdício}\\[1mm]atualiza de hora em hora, alguém olha uma vez por semana};
  \node[caixa, minimum width=46mm, minimum height=26mm, text width=40mm] at (75,43)
    {\textbf{Operacional}\\[1mm]a decisão é de minutos ou horas; atualizar rápido vale o custo};
  \node[caixa, minimum width=46mm, minimum height=26mm, text width=40mm] at (25,15)
    {\textbf{Estratégico}\\[1mm]decisão mensal ou semanal; atualização periódica basta};
  \node[caixa, fill=fundo, minimum width=46mm, minimum height=26mm, text width=40mm] at (75,15)
    {\textbf{Risco}\\[1mm]muita gente decide todo dia com um número velho};

  % eixos
  \draw[seta] (-3,-1) -- (104,-1);
  \draw[seta] (-1,-3) -- (-1,61);
  \node[rotulo, anchor=north] at (50,-2) {com que frequência alguém olha os números};
  \node[rotulo, anchor=south, rotate=90] at (-2,29) {com que frequência os dados atualizam};
  \node[rotulo, anchor=north west] at (0,-2) {menos};
  \node[rotulo, anchor=north east] at (102,-2) {mais};
\end{tikzpicture}

  \caption{Frequência de atualização contra frequência de consumo. Nos quadrantes em
    branco as duas estão alinhadas; nos quadrantes em cinza, não.}
  \label{fig:frequencia-consumo}
  \fonte{elaborado pelo autor.}
\end{figure}

Seja qual for a frequência escolhida, o painel deve mostrar, de forma visível, a data e a hora
da última atualização. O usuário precisa saber a idade do número que está vendo, e um painel
que atualiza todo dia às 7 horas, mas que naquele dia falhou, deveria dizer isso em vez de
mostrar silenciosamente os números de ontem.

\rascunho{
  Relato de abertura, da minha experiência como gerente: um pedido de \enquote{tempo real}
  que, olhando o uso, era uma reunião semanal, e o que foi feito a partir disso.
}

\section{O que as fontes permitem}
\label{sec:limites-fontes}

Se o ritmo da decisão define a frequência necessária, as fontes definem a frequência possível.
Alguns sistemas só permitem extrações em janelas noturnas, para não pesar na operação. Algumas
APIs limitam o número de chamadas por hora. Alguns dados só ficam corretos depois de um
fechamento: o lançamento contábil do dia pode ser ajustado até o fim do mês. E nem todo sistema
oferece captura de mudanças, que permite receber só o que mudou em vez de reler tudo.

Atualizar mais rápido que a fonte muda não traz dado novo, só custo. Se o sistema de origem
consolida as vendas uma vez por dia, rodar o pipeline de hora em hora vai copiar o mesmo dado
vinte e quatro vezes.

Também vale lembrar que o custo e a complexidade não crescem de forma linear com a frequência.
Passar de diário para horário costuma ser simples; passar de horário para contínuo, em
fluxo, exige outra arquitetura, outras ferramentas e outra forma de monitorar
\parencite{reis2022}. E a forma de carregar pesa tanto quanto a frequência: uma carga
incremental, que processa só os registros novos ou alterados, pode tornar viável uma
frequência que seria cara demais com uma carga completa.

\section{O que cada profissional precisa olhar}
\label{sec:papeis-frequencia}

A frequência não é responsabilidade de um papel só. O gerente decide, mas cada papel enxerga
uma parte da conta, e a decisão é melhor quando todas as partes chegam até ele. A
Tabela~\ref{tab:papeis-frequencia} resume o que cada um precisa olhar.

\begin{table}[htbp]
  \centering
  \caption{A frequência de atualização vista por cada papel.}
  \label{tab:papeis-frequencia}
  \small
  \begin{tabularx}{\textwidth}{@{}>{\raggedright\arraybackslash}p{0.27\textwidth}
      >{\raggedright\arraybackslash}X@{}}
    \toprule
    Papel & Com o que precisa se preocupar \\
    \midrule
    Gerente ou líder de dados & Entender a necessidade real do negócio, negociar a
      frequência com o stakeholder, equilibrar custo e valor, acompanhar as métricas de
      uso e revisar a decisão de tempos em tempos. \\
    \addlinespace
    Stakeholder (dono do negócio) & Dizer quando e como decide com o número, e aceitar
      que atualizar mais rápido tem custo. \\
    \addlinespace
    Analista ou desenvolvedor de BI & Desenhar o produto para o ritmo de quem consome,
      mostrar a data da última atualização e não prometer uma frequência que a fonte não
      sustenta. \\
    \addlinespace
    Analytics engineer & Escolher a materialização certa (visão, tabela, incremental)
      para a frequência combinada e testar a atualização do dado. \\
    \addlinespace
    Engenheiro de dados & Conhecer os limites das fontes, o custo de cada execução e
      monitorar atrasos e falhas de carga. \\
    \bottomrule
  \end{tabularx}
  \fonte{elaborado pelo autor.}
\end{table}

O papel do gerente merece um destaque. É ele quem precisa traduzir o pedido do stakeholder em
necessidade real, e quem precisa sustentar a conversa quando a resposta for \enquote{diário é
suficiente}. Isso exige conhecer o negócio bem o bastante para propor a frequência certa, e não
só aceitar a pedida (capítulo~\ref{cap:negocio}). E exige números: a conversa fica muito mais
fácil quando o gerente pode mostrar quantas vezes o painel foi aberto no último mês e quanto
custou cada atualização.

Uma forma prática de formalizar o combinado é o objetivo de nível de serviço
(\veja{sec:observabilidade}): \enquote{os dados do dia anterior estarão disponíveis até as 8
horas em 99\% dos dias úteis}. Ele explicita a frequência e o prazo, e cria uma base objetiva
para monitorar e para rediscutir \parencite{beyer2016}.

\section{Métricas sobre os próprios produtos de dados}
\label{sec:metricas-uso}

Para saber se a frequência combinada está certa, é preciso tratar os painéis, tabelas e
modelos como produtos, com métricas próprias. Algumas das mais úteis:

\begin{itemize}
  \item número de acessos e de usuários distintos por período;
  \item data do último acesso;
  \item número de atualizações por período;
  \item custo de cada atualização e custo total por produto;
  \item custo por acesso, ou atualizações por acesso.
\end{itemize}

A última é a que mais revela desperdício. Um painel que atualiza vinte e quatro vezes por dia
e é aberto duas vezes por semana tem mais de oitenta atualizações por acesso. Quase toda
ferramenta de BI e quase todo armazém de dados guarda, de alguma forma, os registros de quem
acessou o quê e de quanto custou cada execução. A consulta abaixo mostra a ideia, sobre duas
tabelas de metadados hipotéticas.

\begin{sql}
-- Atualizações contra acessos de cada produto de dados nos últimos 30 dias
WITH atualizacoes AS (
    SELECT produto,
           COUNT(*)            AS qtd_atualizacoes,
           SUM(custo_creditos) AS creditos
    FROM meta_atualizacoes
    WHERE executado_em >= CURRENT_DATE - INTERVAL '30' DAY
    GROUP BY produto
),
acessos AS (
    SELECT produto,
           COUNT(*)                AS qtd_acessos,
           COUNT(DISTINCT usuario) AS usuarios
    FROM meta_acessos
    WHERE acessado_em >= CURRENT_DATE - INTERVAL '30' DAY
    GROUP BY produto
)
SELECT a.produto,
       a.qtd_atualizacoes,
       COALESCE(v.qtd_acessos, 0)            AS qtd_acessos,
       COALESCE(v.usuarios, 0)               AS usuarios,
       a.creditos,
       a.creditos / NULLIF(v.qtd_acessos, 0) AS creditos_por_acesso
FROM atualizacoes AS a
LEFT JOIN acessos AS v ON v.produto = a.produto
ORDER BY creditos_por_acesso DESC NULLS FIRST;  -- sem acesso nenhum aparece primeiro
\end{sql}

O resultado é uma lista priorizada. No topo ficam os produtos que consomem recurso e ninguém
abre, candidatos a desligar. Logo abaixo, os que atualizam muito mais do que são consultados,
candidatos a baixar a frequência. A mesma lógica vale para assistentes de IA: tokens gastos por
resposta efetivamente usada são o equivalente, na era dos modelos de linguagem, dos créditos
por acesso.

\section{Revisar a definição}
\label{sec:revisar-frequencia}

A frequência combinada no começo do projeto é uma hipótese. O stakeholder imaginava que
olharia o painel todo dia; a rotina dele mostrou que olha uma vez por semana. Ou o contrário:
um painel pensado para a reunião mensal virou ferramenta diária da operação, e agora atualiza
devagar demais. As métricas de uso dizem qual era a realidade, e a definição precisa ser
revisada com base nelas.

Um ritual simples resolve: uma vez por mês ou por trimestre, revisar os produtos mais caros e
os menos acessados, com a lista da seção anterior na mão. Essa revisão cabe naturalmente na
cadência do Scrum com Kanban (capítulo~\ref{cap:scrumban}), como um item fixo da revisão
periódica.

Baixar a frequência ou desligar um produto também é entrega, e precisa ser tratado como tal.
Isso significa conversar com o dono do produto, mostrar os números de uso, combinar a mudança
e comunicá-la a quem usa. Desligar sem avisar destrói confiança; não desligar nunca destrói
orçamento. Numa época em que quase tudo é cobrado por uso, em créditos de processamento ou em
tokens, o custo de errar nessa definição deixou de ser invisível, e é o assunto do próximo
capítulo.

\begin{destaque}
Atualize o mais rápido que as fontes permitem, mas só o quanto a decisão precisa. E meça o
uso, porque a primeira definição quase sempre está errada.
\end{destaque}

\rascunho{
  Relato: uma revisão de frequência ou de produtos sem uso, e o que ela economizou. Números
  só em ordem de grandeza.
}

\begin{leituras}
  \leitura{beyer2016}{De onde vem a ideia de objetivo de nível de serviço (SLO). Os
    capítulos sobre SLOs se aplicam quase sem adaptação a prazos de atualização de dados.}
  \leitura{moses2022}{Trata a atualização do dado como uma das dimensões da
    observabilidade e mostra como monitorar atrasos antes do usuário reclamar.}
  \leitura{reis2022}{Para a discussão de lote contra fluxo contínuo e do custo de cada
    escolha no ciclo de vida da engenharia de dados.}
\end{leituras}
